\subsection{Third Letter}

\begin{quotex}
This letter was written at the end of 1951, the year of Guenon's death. It is clear that the several letters between Evola and Guenon in the preceding months had changed little in the disagreements between them.

At the beginning of the letter, Evola is referring to Eliade's response to Evola's previous question about Eliade's lack of references to Evola or Guenon in his books. At his post at the Sorbonne, Eliade seems to be hiding his past association with esoterism and right wing movements. Evola is concerned about the negative psychic forces in this choice, but assumes Eliade is still involved in some groups that would prevent that. In contrast, Evola was invited to teach courses on race at two Italian universities, although Evola did not like academic life. If Eliade was truly acting as a Trojan horse to inject certain ideas into Academia, its effects have been nil.

After 25 years, we see Evola is still attached to his system of Absolute Idealism. That system is not fully compatible with Guenon's metaphysical system derived from the Vedanta and other Traditional sources. Those differences affect their respective worldviews in significant ways. Both Evola and Guenon want to ``move beyond" philosophy; for Guenon, that means the Intellect, but for Evola, that is the Will. I should remind readers that 20 years earlier, the young Eliade had written an unpublished commentary on Evola's philosophical system of Absolute Idealism.

Evola's closing to the letter is interesting and out of character. Perhaps it was related to the time of year the letter was written … the Unconquerned Sun was celebrated on Christmas Day.

\end{quotex}
31 Dec 1951

Many thanks for your courteous letter and thanks also for having arranged the sending of your new book [\emph{Shamanism}], that I will read with particular interest. Then I will tell you what are the possibilities at Laterza [a publisher]; it stands to reason that I will do my best to be useful.

Regarding your clarifications about your relations with academic ``masonry", I find them somewhat satisfactory. It would therefore be less a question about methodology than pure tactic, and there would be nothing to say against the attempt to introduce any Trojan horse into the university citadel. The important thing would be to not let yourself take part, in any way, in a deception, because a sort of ``psychic current" meets in academic circles, with the possibility of a subtle deforming and contaminating influence. But I think that, both as the interior foundation, and through your probable relations with circles qualified in a different way, you can defend yourself from this danger.

As to ``methodology", you well know that I seek to follow a middle way since, differently from most ``esoterists", I am also concerned to produce research satisfactory from the ``scientific" point of view. What you undertake in the fields of the science of religions and mythology, I undertook many years ago, but in the field of academic philosophy that was then absolute idealism. The direction was the same: to show that the most important problems of this philosophy cannot be resolved, if it does not go beyond ``philosophy" \emph{tout court}. But after this contribution, expressed in three books (I recently revised one, the \emph{Theory of the Absolute Individual}, and I think it useful for have it republished as an account of it), I had enough of it. I don't know the environment of the Sorbonne. As it concerns Italy, at least until recently it was not necessary to disguise oneself too much, since I myself had received the assignment to teach some courses in the universities of Milan and Florence. But my conclusion was that the game is not worth the candle; and the repulsion for the types and the cabals of the university world is for me physiological.

Since you mention Mr. Guenon in particular, I think that a useful action would consist in developing certain aspects of his doctrine that suffer from a fundamentally arbitrary dogmatism since, all things considered, the mixing of traditional data with individual points of view was inevitable even in his case. So much more in France, but also in Italy, groups were formed that follow the master as though the ``head of the class", redoubling the dogmatic certainty and claiming to be the only ones to administer ``orthodoxy"; that thing is somewhat tiresome and can only be harmful to what is best in Guenon.

I am very obligated to you for your intention to help me get some of my books published in French. With Gallimard and De Noel the thing is only interrupted; as for the latter publisher, the person who mediated and had already begun the translation of \emph{Revolt} has vanished. Regarding the former, after the attempt he made at Leterza, no one any longer knows what happened and although I had written myself, they gave no response. In any case, I think that Payot has some book series in which the two books you mention (\emph{Revolt against the Modern World} and \emph{Doctrine of Awakening}) would fit in rather well, since they would be among the works that are certainly no more ``scientific" than mine. The important thing is that they do not encounter prejudices of principle; these, moreover, could be reduced through the fact that the German and English translations appeared in very ``respectable" publishers. In any case, I think that the best thing is to wait for your return to Paris, before making any attempts and before I send the books. Consequently, I ask you to write me a few words when you return to Paris.

It pleases me to learn that you will come again to Italy. Since I would not want to lose the opportunity for a meeting, I ask you to alert me when your plans are definitive, so I can adjust my schedule, since it is possible that I will leave Rome in the spring for a certain period.

With my best wishes – \emph{in signo solis invicti} – for the new cycle, very cordially yours …



\flrightit{Posted on 2012-10-01 by Cologero }

\begin{center}* * *\end{center}

\begin{footnotesize}\begin{sffamily}



\texttt{Avery on 2012-10-01 at 11:20 said: }

FWIW:

Doctrine of Awakening was published in 1948 by the orientalist publisher ``Luzac \& Co." in London.

Revolt was first published in German by the mainstream house Deutsche Verlags-Anstalt of Stuttgart in 1935.

Revolt against the Modern World was first published in English by IT.

Doctrine has never been published in German: \url{http://www.oocities.org/capitolhill/6890/evolabib.html}.


\hfill

\texttt{cosmodromium on 2012-10-01 at 11:41 said: }

Excellent! Thank You. My Bollingen English trans. of Eliade's YOGA (1958, originally 1954 in French) includes Evola in the index with 5 entries, all footnotes, to Evola's Doctrine of Awakening; Yoga of Power; and The Hermetic Tradition. With the exception of the first, the notes are concentrated in Eliade's chapter on Tantra.


\hfill

\texttt{Izak on 2012-10-02 at 18:23 said: }

Thanks for the translation.

Evola's point about wanting to convey tradition through a means which is valid from the aspect of the ``sciences" is pretty interesting.

I don't know a lick of Italian, but I'm assuming that the language distinguishes between ``sciences" (as just general knowledge disciplines) and ``natural sciences" the same way that German does, which is why Evola then goes on to discuss his philosophy.

But does Evola's claim mean that he is still deeply attached to his brand of magical idealism, or that he just agrees with its very basic purpose, from a more superficial stance?

It seems to me that the chief distinction here is that Evola wants to find a way to translate the language of tradition into a more contemporary language of disciplined discourse, so the real question is one more of language than content.


\hfill

\texttt{Cologero on 2012-10-02 at 19:19 said: }

It's possible to a certain extent, Izak. Yes, by ``scientific", he means in the university sense of an organized body of objective knowledge. That is how he identifies Eliade's approach, although he knows Eliade is holding back. Obviously, Evola's system is not strictly speaking mere philosophy, because it requires a certain spiritual attainment to understand. Keep in mind, too, as we have documented previously, the young Eliade wrote a never-published commentary on Evola's magical idealism.

It is also the case that Evola developed his philosophical system before his encounter with Guenon and tradition. Italy was the center of philosophical idealism at the time and Evola had done an intense study of German and Italian idealism. Now I don't know what the original system was because Evola had revised it for later publications. I have the final edition which does, indeed, contain references to more traditional persons who are not typically found in philosophical treatises. So perhaps, over time, he sought to create a Western philosophical system compatible with Traditional metaphysics.

I agree that is something that needs to be done. I could point out that from a personal point of view, based on my background and education, I am actually closer to Evola's approach. Guenon's metaphysics often seems alien; however, it is necessary to meditate on it to make it come alive.

If there is sufficient interest, I will try to extract relevant passages from Evola's philosophical works for translation. As you can imagine, that is not a trivial task. In the meantime, our translation and commentaries on his The Individual and the Becoming of the World will provide an entry point into his thought.


\hfill

\texttt{Izak on 2012-10-02 at 20:43 said: }

Yes, if you're up for the challenge of translating Evola's more idealist-oriented stuff, that would be great! 

My only exposure to his ``idealist" period has been in his autobiography, and in a discussion by the scholar Paul Furlong (Avery reviewed that book over at Amerika.org, so I read it shortly afterward). Furlong's depiction of Evola's early days is interesting, and he admirably understands both idealistic philosophy as a whole and Evola's relationship to Gentile and Croce. But he also makes some claims/insinuations which I'm not terribly sure about.

I also wonder about the degree to which his theory of the ``absolute individual" does permeate his later work, and the exact moment at which it begins to break off, if ever. Both The Hermetic Tradition and The Yoga of Power do both seem fairly close to what little I know of his school-philosophy days, that is if his ideas have been portrayed accurately in the scholarship. 

In the meantime, I'll check out The Individual and the Becoming of the World.


\hfill

\texttt{Cologero on 2012-10-02 at 20:57 said: }

In the Gornahoor Forum, Graham has made available our commentaries on IABW. If you do a search here on the Individual and the Becoming of the World, you can find most of it translated. I don't know if the Rune Guild is still selling the pdf version of it.

Although I assert the right to offer my own commentaries, unlike ``scholars", I try to also provide translations. That way readers can judge for themselves.

The books you mentioned were certainly early works. As you see from this letter, as late as 1951, after years of correspondence with Guenon, he still regarded absolute idealism as fundamental to his world view. Shortly before his death, he republished the texts in 1973 after extensive revisions. So, the inference must be that he still considered it valid. The only book that Evola did renounce was ``Pagan Imperialism".


\hfill

\texttt{Izak on 2012-10-02 at 22:32 said: }

Yes, and I've been wondering for a while just what his revisions are like, especially considering that The Yoga of Power was originally called Man as Power, so even the titles have different implications (and I'll add that upon your recommendation, I've read Woodroffe's The World As Power, so I at least have that point of reference). 

If mainstream academics like the guy I mentioned were really worth their salt, they would provide translated editions of these books with commentaries containing extensive footnotes, maybe describing the differences between the first edition and later editions, or pointing out factual errors, providing exegetical notes, noting manuscript changes, or whatever else might be relevant. After all, they're getting paid to do this stuff! But I guess they are put in a strange political position, where they must assume no prior knowledge of his work, and to posture themselves as people who are there to put Evola to rest once and for all, often fairly clumsily, for the ideologically innocent academic community of moral do-gooders. This is all quite funny to me, because in the last year, there've been more scholarly attempts in English to grapple with his work than there ever have been before.

But that guy's view and yours seem to be the same, which is that Evola didn't change too much from his original philosophical position. And the letter seems to suggest as much, too.

Thanks for the info!


\end{sffamily}\end{footnotesize}
